\hypertarget{vlaux6a21ux578b}{%
\chapter{VLA模型}\label{vlaux6a21ux578b}}

\hypertarget{vlaux6a21ux578bux7684ux67b6ux6784ux4e0eux8badux7ec3ux65b9ux6cd5}{%
\section{VLA模型的架构与训练方法}\label{vlaux6a21ux578bux7684ux67b6ux6784ux4e0eux8badux7ec3ux65b9ux6cd5}}

VLA（Vision-Language-Action，视觉-语言-动作）模型基于预训练的视觉-语言模型（VLM）微调而来，是大模型时代具身智能领域的重要范式。它打破了传统机器人``感知-规划-控制''的模块化割裂设计，将语义直接映射为物理世界的动作。

\hypertarget{ux4e00ux6838ux5fc3ux601dux60f3-13}{%
\subsection{一、核心思想}\label{ux4e00ux6838ux5fc3ux601dux60f3-13}}

VLA模型的核心思想是``端到端''（End-to-End）的多模态深度融合。它将预训练的视觉-语言模型（VLM）作为``大脑''，把机器人的动作输出转化为一种特殊的``语言''，从而让大模型能够直接指挥物理硬件。

\hypertarget{ux4e00ux8f93ux5165ux8f93ux51faux6a21ux6001}{%
\subsubsection{（一）输入输出模态}\label{ux4e00ux8f93ux5165ux8f93ux51faux6a21ux6001}}

\begin{enumerate}
\def\labelenumi{\arabic{enumi}.}
\item
  视觉（Vision）：通过视觉编码器（如ViT、SigLIP）将输入图像（如RGB摄像机画面流）切分为Patch，并转化为视觉Token，赋予机器人理解空间关系、物理属性的能力。
\item
  语言（Language）：接收人类的自然语言指令，结合预训练大模型积累的庞大互联网知识，提供逻辑推理与``常识''。
\item
  动作（Action）：VLA将动作词表化，将连续的机器人物理控制参数（如机械臂的关节扭矩、末端执行器的坐标）离散化，作为一种特殊的Token融入大语言模型的词表中。
\end{enumerate}

\hypertarget{ux4e8cux57faux672cux67b6ux6784}{%
\subsubsection{（二）基本架构}\label{ux4e8cux57faux672cux67b6ux6784}}

\begin{enumerate}
\def\labelenumi{\arabic{enumi}.}
\item
  VLM：一般取预训练模型进行微调，负责根据视觉和语言输入进行推理和高层决策。
\item
  动作输出头（Action Head）：一般也采用Transformer架构，与VLM组合后共同训练，使得系统由单个VLM转变为能输出机器人动作的VLA。VLM的输出作为上下文输入动作头。
\end{enumerate}

\hypertarget{ux4e8cux52a8ux4f5cux5934ux7528ux81eaux56deux5f52ux7684vlaux8303ux5f0f}{%
\subsection{二、动作头用自回归的VLA范式}\label{ux4e8cux52a8ux4f5cux5934ux7528ux81eaux56deux5f52ux7684vlaux8303ux5f0f}}

\hypertarget{ux4e00ux57faux672cux539fux7406-2}{%
\subsubsection{（一）基本原理}\label{ux4e00ux57faux672cux539fux7406-2}}

假设在时间步 \(t\)，智能体接收到视觉观测序列 \(o_{\le t}\in\mathcal{O}\) 以及自然语言指令 \(l\in\mathcal{L}\)，模型的目标是输出最优动作 \(a_t\in\mathcal{A}\)。

在以 RT-2（Robotic Transformer 2）和 OpenVLA 为代表的自回归（Autoregressive）VLA 范式中，连续动作 \(a_t\) 会被离散化为一系列动作 Token \((z_1,z_2,\ldots,z_K)\)。例如，一个 7 自由度（7-DoF）的机械臂动作可以表示为空间位移、旋转和夹爪开合度。

模型的训练目标是最大化以观测和指令为条件的动作 Token 的对数似然（Log-Likelihood）：

\[
\begin{aligned}
\mathcal{J}(\theta)
&=
\mathbb{E}_{(o,l,a)\sim\mathcal{D}}
\left[
\sum_{k=1}^{K}
\log P_{\theta}(z_k\mid z_{<k}, o_{\le t}, l)
\right]
\end{aligned}
\]

其中，\(\theta\) 为神经网络的参数，\(\mathcal{D}\) 为包含人类专家示范和互联网图文的混合数据集。通过这种方式，动作预测被完全等价转化为了类似于 ChatGPT 生成文本的 Next-token Prediction 任务。

\hypertarget{ux4e8cux5de5ux4f5cux6d41-4}{%
\subsubsection{（二）工作流}\label{ux4e8cux5de5ux4f5cux6d41-4}}

\begin{enumerate}
\def\labelenumi{\arabic{enumi}.}
\item
  多模态输入处理（Multimodal Input Processing）：步骤 1.1，获取当前相机的 RGB 图像 \(o_t\)，通过视觉编码器（Vision Encoder）提取特征，并经过 Projector（投影层）对齐到语言模型的特征空间，生成视觉 Token 序列；步骤 1.2，接收用户指令 \(l\)，通过文本分词器生成语言 Token 序列。
\item
  特征拼接与跨模态注意力计算（Context Concatenation \& Cross-modal Attention）：步骤 2.1，将系统提示词、视觉 Token 和语言 Token 按序列拼接为 \texttt{{[}System\ Prompt,\ Image\ Tokens,\ Language\ Tokens{]}}；步骤 2.2，送入 Transformer Decoder 骨干网络，通过自注意力机制（Self-Attention）让指令语义与图像中的物理特征产生深度绑定（Grounding）。
\item
  自回归动作生成（Autoregressive Action Generation）：步骤 3.1，Transformer 结合上下文，逐个预测当前步骤的动作 Token。例如先预测代表 \(A_x\) 的 Token，再将 \(A_x\) 作为已知条件预测 \(A_y\)，依次类推，直到输出夹爪动作 Token。
\item
  反词表化与物理控制（Detokenization and Execution）：步骤 4.1，将生成的动作 Token 序列 \((z_1,\ldots,z_K)\) 映射回连续的物理数值，生成实际控制指令；步骤 4.2，将指令发送给机器人底层控制器（如 PID 控制器）驱动电机执行。
\end{enumerate}

\hypertarget{ux4e09ux4f18ux7f3aux70b9}{%
\subsubsection{（三）优缺点}\label{ux4e09ux4f18ux7f3aux70b9}}

自回归VLA直接复用LLM的架构，天然继承了互联网级数据的逻辑链条（Chain-of-Thought）和常识推理能力。泛化能力极强，但受限于自回归的生成速度，高频连续控制能力略显不足。

\hypertarget{ux4e09ux52a8ux4f5cux5934ux7528ux6269ux6563ux6d41ux5339ux914dux7684vlaux8303ux5f0f}{%
\subsection{三、动作头用扩散/流匹配的VLA范式}\label{ux4e09ux52a8ux4f5cux5934ux7528ux6269ux6563ux6d41ux5339ux914dux7684vlaux8303ux5f0f}}

这一范式的核心动机是消除``动作词表化''带来的离散化误差。它不再将动作强行转化为语言词表中的 Token，而是使用生成式连续时间模型（Diffusion Models 或 Flow Matching）直接在连续空间中生成高维、高频的动作轨迹（Action Chunks）。代表性模型包括 Physical Intelligence 的 \(\pi_0\)（Pi-Zero）、ForceVLA 等。

扩散/流匹配 VLA 通常采用 \textbf{``VLM + 动作专家（Action Expert）''} 的解耦与重组架构：

\begin{enumerate}
\def\labelenumi{\arabic{enumi}.}
\item
  多模态条件编码器：继承预训练的视觉语言模型（如 PaliGemma、SigLIP 结合 LLM 骨干），负责将图像序列 \(o\) 和语言指令 \(l\) 处理为深度的上下文条件嵌入（Condition Embeddings）。
\item
  连续动作生成器（Action Expert）：通常是 Diffusion Transformer（DiT）或 U-Net，接收上述条件嵌入，并通过交叉注意力机制（Cross-Attention）或拼接自注意力机制，对纯噪声信号进行迭代去噪，最终输出连续的物理动作块 \(A_t=(a_t,a_{t+1},\ldots,a_{t+H-1})\)。
\end{enumerate}

以目前先进的最优传输流匹配（Optimal Transport Flow Matching，类似于 \(\pi_0\) 中的设计）为例。假设基础噪声分布为 \(x_0\sim p_0=\mathcal{N}(0,I)\)，真实动作轨迹分布为 \(x_1\sim p_1\)，多模态条件为 \(c=\mathrm{Encoder}(o,l)\)。模型构建一条连接噪声与真实动作的直线概率路径：

\[
x_t=(1-t)x_0+t x_1
\]

动作生成器 \(v_\theta\) 的目标是拟合真实的速度场（Velocity Field），其目标函数定义为：

\[
\begin{aligned}
\mathcal{L}_{\mathrm{FM}}(\theta)
&=
\mathbb{E}_{\substack{t\sim \mathcal{U}(0,1)\\ x_0\sim p_0,\ x_1\sim p_1}}
\left[
\left\|v_\theta(x_t,t,c)-(x_1-x_0)\right\|^2
\right]
\end{aligned}
\]

这里 \((x_1-x_0)\) 是真实路径的恒定梯度。相较于传统的 DDPM 扩散模型，流匹配的 ODE 轨迹更加平滑且笔直，能够在极少的推理步数（甚至几步内）生成高保真的连续动作信号。
工作流：

\begin{enumerate}
\def\labelenumi{\arabic{enumi}.}
\item
  特征提取：VLM 接收当前相机观测和语言 \(l\)，前向传播生成多模态条件向量 \(c\)。
\item
  噪声初始化：从标准正态分布中采样一段长度为 \(H\) 的纯噪声轨迹 \(x_0\sim\mathcal{N}(0,I)\)。
\item
  常微分方程求解（ODE Solving）：在时间步 \(t\in[0,1]\) 之间，动作生成器 \(v_\theta(x_t,t,c)\) 结合上下文 \(c\) 计算当前状态的梯度，并通过数值积分器（如 Euler 法或 Heun 法）逐步更新。
\item
  动作执行：将 ODE 终点作为预测的动作块。系统采用模型预测控制（MPC/Receding Horizon Control）策略，仅执行前几个动作，随后重新获取新的视觉观测，进入下一轮闭环推理。
  优缺点：动作输出不再是离散 Token，而是通过 Diffusion Model（扩散模型）或 Flow Matching 直接生成连续的动作轨迹，在处理需要极高精度和灵巧度（如折叠衣服、灵巧手操作）的任务时表现出显著优势，但在长程任务的因果关系保持上可能不如自回归。
\end{enumerate}

\hypertarget{ux56dbux7edfux4e00ux8868ux5f81ux4e0eux6f5cux5728ux52a8ux4f5cux8303ux5f0f}{%
\subsection{四、统一表征与潜在动作范式}\label{ux56dbux7edfux4e00ux8868ux5f81ux4e0eux6f5cux5728ux52a8ux4f5cux8303ux5f0f}}

不同机器人的硬件形态之间（乃至与人类第一视角之间）存在差异。为了让VLA模型不被限制于特定的硬件中，而是可以利用海量无动作标签的视频进行扩展（Scaling），我们引入``抓取右上方的物体''等与具体硬件无关的统一潜在动作，而不是``右臂第一个关节上抬30厘米''这样的具体机械操作。

在具身智能架构分类中，它介于分层和端到端之间，因为模型间传递的信息是可微的，梯度可以（在某些设计下）从底层直接传导回高层，或者在同一个世界模型的潜空间内完成对齐。

\hypertarget{ux4e00ux6838ux5fc3ux67b6ux6784}{%
\subsubsection{（一）核心架构}\label{ux4e00ux6838ux5fc3ux67b6ux6784}}

该路线的核心架构仍可拆成两部分：多模态条件编码器继承预训练的 VLM，将图像序列 \(o\) 与语言指令 \(l\) 编码为上下文条件嵌入；连续动作生成器接收这些条件嵌入，生成统一潜在动作块 \(A_t=(a_t,a_{t+1},\ldots,a_{t+H-1})\)。区别在于，高层输出的并不一定是某个机器人专属的底层控制量，而是可跨形态共享的潜在动作表示。

\hypertarget{ux4e8cux8badux7ec3ux9636ux6bb5ux4e00ux6f5cux5728ux52a8ux4f5cux7a7aux95f4ux7684ux65e0ux76d1ux7763ux5b66ux4e60}{%
\subsubsection{（二）训练阶段一：潜在动作空间的无监督学习}\label{ux4e8cux8badux7ec3ux9636ux6bb5ux4e00ux6f5cux5728ux52a8ux4f5cux7a7aux95f4ux7684ux65e0ux76d1ux7763ux5b66ux4e60}}

这一阶段不需要任何底层物理控制信号（如关节扭矩），模型仅通过观看海量视频序列来学习。模型通常引入一个基于 VQ-VAE（向量量化变分自编码器）的逆动力学模型（Inverse Dynamics Model）：输入连续两帧视频图像 \((o_t,o_{t+1})\)，编码器提取两帧之间的视觉和物理状态变化，生成连续隐特征；向量量化模块将连续隐特征映射到离散密码本（Codebook）中，提取以任务为中心的潜在动作 \(z_t\)；前向动力学模型接收当前帧 \(o_t\) 和潜在动作 \(z_t\)，尝试重建下一帧图像 \(\hat{o}_{t+1}\)。

损失函数由重建损失和密码本正则化损失组成：

\[
\begin{aligned}
\mathcal{L}_{\mathrm{VQ}}
&= \|o_{t+1}-\hat{o}_{t+1}\|^2
+\|\mathrm{sg}[\mathrm{Encoder}(o_t,o_{t+1})]-e_{z_t}\|^2
\\
&\quad+
\beta\|\mathrm{Encoder}(o_t,o_{t+1})-\mathrm{sg}[e_{z_t}]\|^2
\end{aligned}
\]

其中，\(e_{z_t}\) 为密码本中的嵌入向量，\(\mathrm{sg}[\cdot]\) 为停止梯度操作。

\hypertarget{ux4e09ux8badux7ec3ux9636ux6bb5ux4e8cux8054ux5408ux7b56ux7565ux4f18ux5316ux4e0eux5177ux8eabux5bf9ux9f50}{%
\subsubsection{（三）训练阶段二：联合策略优化与具身对齐}\label{ux4e09ux8badux7ec3ux9636ux6bb5ux4e8cux8054ux5408ux7b56ux7565ux4f18ux5316ux4e0eux5177ux8eabux5bf9ux9f50}}

在获得跨形态潜在动作空间后，开始正式训练 VLA 大模型，将语义与物理硬件绑定。首先利用阶段一训练好的 VQ-VAE 编码器，把带有语言指令的机器人视频数据转化为``指令-图像-潜在动作 \(z_t\)''的离散序列数据；然后用标准的 Next-token Prediction 目标训练 VLM 主干网络，使其在给定图像和文本上下文时输出正确的潜在动作分布 \(P_\theta(z_t\mid o_{\le t},l)\)；最后引入轻量级形态感知解码器（Embodiment-aware Decoder）或联合扩散模块，将 VLM 输出的潜在动作 \(z_t\) 与当前机器人的专属形态标识拼接，并使用少量包含真实底层动作标签的数据训练该解码器：

\[
\begin{aligned}
\mathcal{L}_{\mathrm{Action}}
&=
\mathbb{E}\left[
-\log P_\phi(a_t\mid z_t,o_t,\mathrm{embodiment\_id})
\right]
\end{aligned}
\]

其中，目前业界部分实践在高层运用自回归以适配成熟生态，底层运用扩散以降低延迟。

\hypertarget{ux4e94vlaux7684ux7f3aux9677}{%
\subsection{五、VLA的缺陷}\label{ux4e94vlaux7684ux7f3aux9677}}

由于VLA模型本质上由语言模型微调而来，因此它仍然以自然语言空间作为``大脑''的核心表征层，动作则作为一种附属的输出模态。这种设计使得VLA在语义理解和任务泛化方面具有明显优势，但也意味着物理世界的连续状态、动力学约束和动作后果并没有被模型直接建模。

此外，VLA通常依赖当前观测和语言指令直接预测动作，在长时序任务中容易累积误差；当机器人进入训练数据覆盖较少的状态时，模型也缺乏显式的未来预测和自我纠错机制。因此，如何引入世界模型、强化学习或更统一的物理表征，是进一步提升VLA可靠性与长期规划能力的重要方向。

\hypertarget{ux53c2ux8003ux6587ux732e-151}{%
\subsection{参考文献}\label{ux53c2ux8003ux6587ux732e-151}}

\vspace{0.25em}
\begin{itemize}
\tightlist
\item
  Brohan, A., Brown, N., Carbajal, J., et al.~(2023). \href{https://arxiv.org/abs/2307.15818}{RT-2: Vision-Language-Action Models Transfer Web Knowledge to Robotic Control}. arXiv:2307.15818.
\item
  Kim, M. J., Pertsch, K., Karamcheti, S., et al.~(2024). \href{https://arxiv.org/abs/2406.09246}{OpenVLA: An Open-Source Vision-Language-Action Model}. arXiv:2406.09246.
\item
  Black, K., Brown, N., Driess, D., et al.~(2024). \href{https://arxiv.org/abs/2410.24164}{π0: A Vision-Language-Action Flow Model for General Robot Control}. arXiv:2410.24164.
\item
  Ye, S., Jang, J., Jeon, B., et al.~(2024). \href{https://arxiv.org/abs/2410.11758}{Latent Action Pretraining from Videos}. arXiv:2410.11758.
\item
  Bu, Q., Yang, Y., Cai, J., et al.~(2025). \href{https://arxiv.org/abs/2505.06111}{UniVLA: Learning to Act Anywhere with Task-centric Latent Actions}. arXiv:2505.06111.
\item
  Chen, Y., Ge, Y., Tang, W., et al.~(2024). \href{https://arxiv.org/abs/2412.04445}{Moto: Latent Motion Token as the Bridging Language for Learning Robot Manipulation from Videos}. arXiv:2412.04445.
\end{itemize}

\clearpage

\hypertarget{vlaux6a21ux578bux7684ux5f3aux5316ux5b66ux4e60}{%
\section{VLA模型的强化学习}\label{vlaux6a21ux578bux7684ux5f3aux5316ux5b66ux4e60}}

\hypertarget{ux4e00vlaux6a21ux578bux5f3aux5316ux5b66ux4e60ux7684ux96beux70b9}{%
\subsection{一、VLA模型强化学习的难点}\label{ux4e00vlaux6a21ux578bux5f3aux5316ux5b66ux4e60ux7684ux96beux70b9}}

普通连续控制策略一般写作：

\[
a_t\sim\pi_\theta(a_t\mid s_t)
\]

以高斯策略为例：

\[
\pi_\theta(a\mid s)=\mathcal{N}\bigl(\mu_\theta(s),\sigma_\theta(s)\bigr)
\]

于是SAC、PPO等算法可以直接计算 \(\log\pi_\theta(a_t\mid s_t)\)，并据此得到策略梯度项：

\[
\nabla_\theta\log\pi_\theta(a_t\mid s_t)A_t
\]

对于自回归VLA，策略是每一步概率的积，对数策略是每一步对数概率的和；但扩散/流匹配VLA就难以计算对数策略了。而且如果让RL梯度穿过完整的10-50步diffusion/flow trajectory，不仅显存和计算量大，还很容易不稳定。

所以VLA RL一个很重要的研究问题其实就是：到底应该在哪里做RL？是整个VLA？action head？action token？flow denoising trajectory？还是latent noise？

\hypertarget{ux4e8cvla-rlux76f4ux63a5ux5bf9vlaux6574ux4f53ux7528ppoux4f18ux5316}{%
\subsection{二、VLA-RL：直接对VLA整体用PPO优化}\label{ux4e8cvla-rlux76f4ux63a5ux5bf9vlaux6574ux4f53ux7528ppoux4f18ux5316}}

\hypertarget{ux4e00ux6838ux5fc3ux601dux60f3-14}{%
\subsubsection{（一）核心思想}\label{ux4e00ux6838ux5fc3ux601dux60f3-14}}

它的想法跟LLM的RLHF/PPO非常接近。OpenVLA原本做的是：

\[
(\text{image},\text{instruction})
\longrightarrow
[v_1,v_2,\ldots,v_7]
\longrightarrow
a_t
\]

其中7个Token对应7个机器人动作维度。因为它采用自回归方式：

\[
\pi_\theta(a_t\mid o_t,l)=\prod_{i=1}^{7}p_\theta(v_i\mid v_{<i},o_t,l)
\]

所以动作的对数概率可以直接得到：

\[
\log\pi_\theta(a_t\mid o_t,l)=\sum_{i=1}^{7}\log p_\theta(v_i\mid v_{<i},o_t,l)
\]

对于OpenVLA这类自回归动作分词器，动作概率很容易计算。论文只需把单步机器人的多模态观测与动作看成类似LLM的多轮对话，然后直接使用PPO与GAE更新VLA。

运用标准PPO算法，对整个VLA模型进行训练即可。

\hypertarget{ux4e8cux7a00ux758fux53cdux9988ux95eeux9898ux53caux5176ux89e3ux51b3}{%
\subsubsection{（二）稀疏反馈问题及其解决}\label{ux4e8cux7a00ux758fux53cdux9988ux95eeux9898ux53caux5176ux89e3ux51b3}}

机器人任务反馈稀疏，比如任务：把杯子放进柜子。可能执行100步后，环境才反馈结果，那么前面第20步``机械臂已经正确抓住杯子''到底好不好，单靠terminal reward很难判断。VLA-RL运用Process Reward Model，根据夹爪开闭的显著变化划分milestone/subtask，并结合末端执行器速度接近零等信息找关键进展状态，从而给中间动作提供过程奖励。

\hypertarget{ux4e09ux3c0rlux8ba9ux6d41ux5339ux914dvlaux6a21ux578bux4e5fux53efux6574ux4f53ux7528rlux8badux7ec3}{%
\subsection{三、πRL：让流匹配VLA模型也可整体用RL训练}\label{ux4e09ux3c0rlux8ba9ux6d41ux5339ux914dvlaux6a21ux578bux4e5fux53efux6574ux4f53ux7528rlux8badux7ec3}}

πRL 中提出的Flow-Noise等方法在原有流匹配VLA模型的基础上引入可学习的噪声网络，将去噪过程重塑为离散时间的MDP，从而可用传统RL算法，既解决了流匹配缺失随机性、无法计算动作概率似然的问题，又通过可学习、有方向的探索解决了高维空间盲目探索往往失败的问题。

\hypertarget{ux4e00ux6838ux5fc3ux601dux60f3-15}{%
\subsubsection{（一）核心思想}\label{ux4e00ux6838ux5fc3ux601dux60f3-15}}

传统流匹配是在连续时间内求解常微分方程（ODE），缺乏随机性，难以像PPO那样计算对数似然。Flow-Noise的创新在于，它在原本确定性的去噪路径中主动注入一个可学习的噪声网络，将去噪过程重塑为离散时间的马尔可夫决策过程（Discrete-time MDP）。这样，模型在每一步生成时既有探索能力，又能获得精确的动作对数似然，从而完整适配传统Actor-Critic在线强化学习算法。

\hypertarget{ux4e8cux7b97ux6cd5ux539fux7406}{%
\subsubsection{（二）算法原理}\label{ux4e8cux7b97ux6cd5ux539fux7406}}

Flow-Noise放弃了纯确定性的ODE采样，引入可学习的噪声网络（Learnable Noise Network），将原本连续的去噪轨迹离散化为若干步，并在每一步注入受控随机性。

在该内部去噪MDP中：

\begin{itemize}
\tightlist
\item
  状态空间：在去噪时间步 \(\tau\)，内部状态定义为当前环境观测与当前噪声动作块的组合：
\end{itemize}

\[
\bar{s}_t^\tau=(o_t,A_t^\tau).
\]

\begin{itemize}
\tightlist
\item
  动作空间：内部``动作''是为流模型生成下一离散时间步隐状态。当 \(\tau<1\) 时，策略输出：
\end{itemize}

\[
\bar{a}_t^\tau=A_t^{\tau+\delta}.
\]

当 \(\tau=1\) 时，输出最终动作 \(A_t^1\)。

\begin{itemize}
\tightlist
\item
  随机转移与可学习噪声：从 \(A_t^\tau\) 到 \(A_t^{\tau+\delta}\) 的转移被建模为高斯分布：
\end{itemize}

\[
A_t^{\tau+\delta}=\mu_\tau+\sigma_\tau\delta\cdot\varepsilon,
\qquad \varepsilon\sim\mathcal{N}(0,I).
\]

其中 \(\mu_\tau\) 由流匹配网络 \(v_\theta\) 预测的均值方向决定， \(\sigma_\tau\) 则是可学习噪声网络输出的标准差。

该设计带来两点效果：

\begin{itemize}
\tightlist
\item
  自适应探索：在生成轨迹的安全区域，网络会输出较大的方差以鼓励广泛探索；当进入对噪声敏感的高频细节生成区域时， \(\Sigma_\phi\) 会自动收敛，输出较小方差，甚至趋近于0。
\item
  维度解耦：网络可以学习在不重要的背景特征维度上加大探索，而在关键的语义维度上保持克制，从而实现高维空间中的精细探索。
\end{itemize}

由于每一步 \(\tau\to\tau+\delta\) 都服从已知的高斯分布，整条去噪轨迹的联合概率密度可以按链式法则分解：

\[
\log\pi_\theta(a_t\mid o_t)=\sum_{\tau=0}^{1-\delta}\log p_\theta
\bigl(A_t^{\tau+\delta}\mid A_t^\tau,o_t\bigr).
\]

因为 \(p_\theta\) 是由参数化高斯分布构成的，上式可以直接计算，从而可用于PPO等基于策略梯度的强化学习目标。

\hypertarget{ux4e09ux5de5ux4f5cux6d41}{%
\subsubsection{（三）工作流}\label{ux4e09ux5de5ux4f5cux6d41}}

\textbf{步骤1：环境交互与初始化（外部循环起点）}

\begin{enumerate}
\def\labelenumi{\arabic{enumi}.}
\tightlist
\item
  智能体从仿真环境或现实世界获取当前时间步 \(t\) 的观测状态 \(o_t\)，如摄像头图像、本体感受数据。
\item
  在模型内部，初始化流匹配的起点，从标准高斯噪声采样 \(A_t^0\sim\mathcal{N}(0,I)\)。
\end{enumerate}

\textbf{步骤2：离散化去噪过程（内部MDP循环）}

开启一个从 \(\tau=0\) 到 \(\tau=1\)、步长为 \(\delta\) 的离散去噪循环。对于每一步：

\begin{enumerate}
\def\labelenumi{\arabic{enumi}.}
\tightlist
\item
  构建内部状态，将环境观测与当前变量组合为 \(\bar{s}_t^\tau=(o_t,A_t^\tau)\) 。
\item
  网络推理：将内部状态 \(\bar{s}_t^\tau\) 输入VLA骨干网络（如Transformer），预测当前速度场 \(v_\theta\)，进而计算期望去噪方向 \(\mu_\tau\)。
\item
  噪声评估：辅助的可学习噪声网络根据当前状态输出标准差 \(\sigma_\tau\)。
\item
  采样与记录：采样随机噪声 \(\varepsilon\)，计算下一步隐变量 \(A_t^{\tau+\delta}=\mu_\tau+\sigma_\tau\delta\cdot\varepsilon\)；同时，将该步的对数似然 \(\log p_\theta(A_t^{\tau+\delta}\mid A_t^\tau,o_t)\) 记录到缓存中。
\end{enumerate}

\textbf{步骤3：动作执行与奖励获取（外部循环终点）}

\begin{enumerate}
\def\labelenumi{\arabic{enumi}.}
\tightlist
\item
  当内部去噪循环到达 \(\tau=1\) 时，得到无噪声的确定性物理动作 \(a_t=A_t^1\)。
\item
  将动作 \(a_t\) 下发给环境执行。
\item
  观测状态更新，并返回单步奖励 \(r_t\)，例如抓取任务成功时返回稀疏的 \(+1\) 奖励。
\end{enumerate}

\textbf{步骤4：强化学习策略更新（RL反向传播）}

\begin{enumerate}
\def\labelenumi{\arabic{enumi}.}
\tightlist
\item
  收集到一个或多个Episode的完整轨迹后，计算累计优势函数。
\item
  提取步骤2中精确计算的联合对数似然总和 \(\sum\log p_\theta\)。
\item
  应用PPO算法的截断代理目标函数计算策略梯度。
\item
  更新VLA模型网络参数及噪声网络参数，使得在相似观测下，能获得高环境奖励的去噪轨迹生成概率最大化。
\end{enumerate}

\hypertarget{ux56dbux8be5ux65b9ux6cd5ux5982ux4f55ux89e3ux51b3ux6d41ux5339ux914dux8131ux79bbux6700ux4f18ux4f20ux8f93ux540eux6b65ux6570ux6fc0ux589eux7684ux95eeux9898}{%
\subsubsection{（四）该方法如何解决流匹配``脱离最优传输后步数激增''的问题}\label{ux56dbux8be5ux65b9ux6cd5ux5982ux4f55ux89e3ux51b3ux6d41ux5339ux914dux8131ux79bbux6700ux4f18ux4f20ux8f93ux540eux6b65ux6570ux6fc0ux589eux7684ux95eeux9898}}

在传统流匹配中，轨迹弯曲会导致ODE求解器在积分时产生巨大的截断误差，从而必须将步长 \(\Delta t\) 切得很小，即增加函数评估次数（NFE）。Flow-Noise从根本上绕开了ODE求解器：它将整个生成过程重构为一个离散时间MDP，环境转移函数被强行定义为

\[
x_{t-\Delta t}=a_t.
\]

这意味着模型不再预测一条连续的``切线方向''，而是直接跳跃到下一个离散状态。强化学习算法优化的是这一离散跳跃点的累积奖励，既不再需要用连续微分方程求积分，也能绕开弯曲轨迹导致截断误差的问题，模型因而可以在较少步数内得到高奖励结果。

这让Flow-Noise看起来很像扩散模型，但同时又具备了流匹配模型自身去噪轨迹接近直线、速度快的优点。

\hypertarget{ux56dbpolicy-agnostic-rlux8ba9rlux68afux5ea6ux4f18ux5316ux8f93ux51faactionux800cux975evlaux6a21ux578b}{%
\subsection{四、Policy Agnostic RL：让RL梯度优化输出Action而非VLA模型}\label{ux56dbpolicy-agnostic-rlux8ba9rlux68afux5ea6ux4f18ux5316ux8f93ux51faactionux800cux975evlaux6a21ux578b}}

既然直接优化policy很困难，那可以先直接优化action。

当前策略为

\[
a_i\sim\pi_\phi(a_i\mid s).
\]

先生成一组候选动作 \(a_1,a_2,\ldots,a_K\)，然后由Critic \(Q_\theta(s,a_i)\) 给每个动作评分。

对每个候选动作执行梯度上升：

\[
a_i^{j+1}=a_i^j+\alpha\nabla_a Q_\theta(s,a_i^j).
\]

注意这里直接对Action做梯度上升，计算的是对Action而不是对VLA模型参数的梯度。

经过迭代后，可以得到最优动作 \(a^\ast\)，或者得到一组优化后的动作 \(a_1^\ast,a_2^\ast,\ldots\)。

这时只需用监督学习将其蒸馏给VLA：

例如，对OpenVLA可采用负对数似然蒸馏：

\[
\mathcal{L}=-\log\pi_\phi(a^\ast\mid s).
\]

对Diffusion Policy可采用去噪目标蒸馏：

\[
\mathcal{L}=\left\|\varepsilon-\varepsilon_\phi(a_t,s,t)\right\|^2.
\]

RL gradient不直接进入VLA。

\hypertarget{ux4e94dsrlvlaux6839ux672cux4e0dux8badux7ec3ux53eaux8badux7ec3ux9009ux4ec0ux4e48ux566aux58f0}{%
\subsection{五、DSRL：VLA根本不训练，只训练``选什么噪声''}\label{ux4e94dsrlvlaux6839ux672cux4e0dux8badux7ec3ux53eaux8badux7ec3ux9009ux4ec0ux4e48ux566aux58f0}}

\hypertarget{ux4e00ux6838ux5fc3ux601dux60f3-16}{%
\subsubsection{（一）核心思想}\label{ux4e00ux6838ux5fc3ux601dux60f3-16}}

一般情况下，扩散模型在推理时会随机选取噪声。论文《Steering Your Diffusion Policy with Latent Space Reinforcement Learning》提出，可以让RL梯度不训练模型，而是训练选择噪声的策略，在冻结模型的前提下，通过输出不同方向的初始噪声将模型的输出指引向不同的动作模式（如w1代表向左、w2代表抓得更深等）。

对比：

一般方法：随机噪声 → 经RL训练后的模型 → 输出动作

DSRL：经RL优化的噪声 → 冻结的预训练模型 → 输出动作

\hypertarget{ux4e8ccriticux8badux7ec3ux4e2dux7684noise-aliasing}{%
\subsubsection{（二）Critic训练中的noise aliasing}\label{ux4e8ccriticux8badux7ec3ux4e2dux7684noise-aliasing}}

DSRL中，Critic模型Q(s,a)需要输入状态s和策略模型给出的机器人动作a。而冻结的策略模型a=F(s,w)，这里F是冻结的，w是输入的噪声，即a是w的函数，即使某个w从来没在真实机器人上执行过，只要知道它映射到什么action，就可以估计它的价值。故Q(s,a)=Q(s,F(s,w))=Q\_w(s,w)。

DSRL-NA先用机器人数据(s,a,s',r)训练出Q(s,a)，进而就可以训练得到Q\_w(s,w)，方法为：对于每一个w，通过函数a=F(s,w)得到a，让Q\_w(s,w)蒸馏Q(s,a)的输出。

\hypertarget{ux516drldgrlux53eaux8d1fux8d23ux751fux4ea7ux76d1ux7763ux5b66ux4e60ux6837ux672c}{%
\subsection{六、RLDG：RL只负责生产监督学习样本}\label{ux516drldgrlux53eaux8d1fux8d23ux751fux4ea7ux76d1ux7763ux5b66ux4e60ux6837ux672c}}

论文《RLDG: Robotic Generalist Policy Distillation via Reinforcement Learning》提出，可用RL分别训练不同任务的较小``专家模型''，然后让它们各自rollout expert trajectories，用这些数据作为大模型监督学习的样本，不用RL直接微调大模型。

由于RL的策略比没有人类策略中存在的少量犹豫、停顿、晃动等，且rollout轨迹里常存在从错误状态恢复的数据，故该方法得到的样本用于训练的效果在一些情况下可能反超人类数据。

\hypertarget{ux53c2ux8003ux6587ux732e-152}{%
\subsection{参考文献}\label{ux53c2ux8003ux6587ux732e-152}}

\vspace{0.25em}
\begin{itemize}
\tightlist
\item
  Lu, G., Guo, W., Zhang, C., et al.~(2025). \href{https://arxiv.org/abs/2505.18719}{VLA-RL: Towards Masterful and General Robotic Manipulation with Scalable Reinforcement Learning}. arXiv:2505.18719.
\item
  Chen, K., Liu, Z., Zhang, T., et al.~(2025). \href{https://arxiv.org/abs/2510.25889}{πRL: Online RL Fine-tuning for Flow-based Vision-Language-Action Models}. arXiv:2510.25889.
\item
  Mark, M. S., Gao, T., Sampaio, G. G., et al.~(2024). \href{https://arxiv.org/abs/2412.06685}{Policy Agnostic RL: Offline RL and Online RL Fine-Tuning of Any Class and Backbone}. arXiv:2412.06685.
\item
  Wagenmaker, A., Nakamoto, M., Zhang, Y., et al.~(2025). \href{https://arxiv.org/abs/2506.15799}{Steering Your Diffusion Policy with Latent Space Reinforcement Learning}. arXiv:2506.15799.
\item
  Xu, C., Li, Q., Luo, J., \& Levine, S. (2024). \href{https://arxiv.org/abs/2412.09858}{RLDG: Robotic Generalist Policy Distillation via Reinforcement Learning}. arXiv:2412.09858.
\end{itemize}

\clearpage
